\begin{explanationbox}
	\textbf{\textcolor{CExplanation}{一句话直觉}}

	托勒密定理的本质是：通过旋转相似，将两条对角线的长度关系转化为两组对边长度乘积之和。

	\medskip
	\textbf{\textcolor{CExplanation}{核心拆解}}
	\begin{description}[style=nextline, leftmargin=2.8em, labelsep=0.5em, itemsep=0.45em]
		\item[\textbf{要点一（旋转相似）：}] 在圆内接四边形中，利用圆周角相等构造两个相似三角形，将对角线分割为两段，每段与一组对边对应。
		\item[\textbf{要点二（乘积相加）：}] 由相似得到比例式，交叉相乘后得到 $AB\cdot CD = AC\cdot BE$ 和 $AD\cdot BC = AC\cdot ED$，相加后 $AC\cdot(BE+ED)=AC\cdot BD$。
		\item[\textbf{要点三（共圆前提）：}] 四点共圆保证了圆周角相等，这是构造相似的基础。
	\end{description}

	\medskip
	\textbf{\textcolor{CExplanation}{几何本质（旋转缩放）}}

	圆内接四边形中，通过旋转三角形，可以将对角线转化为两条对边的线性组合，最终表示为两个乘积的和。

	\medskip
	\textbf{\textcolor{CExplanation}{代数意义（方程关系）}}

	定理给出了六条线段之间的一个二次关系：$ef=ac+bd$。已知其中五个量，可解第六个；已知四边可直接求对角线乘积，无需角度信息。

	\medskip
	\textbf{\textcolor{CExplanation}{考点价值（求值判定）}}
	\begin{itemize}[leftmargin=*, itemsep=0.5em, topsep=0.4em]
		\item \textbf{考法一（求对角线长）：} 给出四边形边长或部分边长与一条对角线，利用定理求另一条对角线。
		\item \textbf{考法二（判断共圆）：} 给出六条线段长度，验证等式是否成立，若成立则四点共圆。
	\end{itemize}

	\medskip
	\textbf{\textcolor{CExplanation}{顿悟点}}

	将一条对角线视为被分成的两段之和，每段与一组对边构成相似三角形的对应边，从而把几何量关系转化为代数等式。

	\medskip
	\textbf{\textcolor{CExplanation}{使用场景}}
	\begin{itemize}[leftmargin=*, itemsep=0.5em, topsep=0.4em]
		\item \textbf{场景一（求边或对角）：} 已知四边求对角线乘积，或已知三边及一条对角线求另一边。
		\item \textbf{场景二（共圆判定）：} 在四边形的边长和对角线已知或可测时，用定理判定四点共圆。
	\end{itemize}
\end{explanationbox}
